
\label{sec:conclusion}

Las conclusiones se organizan por objetivo. Cada una se apoya en una cifra de la Sección~\ref{sec:resultados} y su alcance se discute en la Sección~\ref{sec:discusion}.

%--------------------------------------------------------------------------------------
\section{Conclusión asociada al objetivo general}

El objetivo general pedía una metodología para estimar la tasa de desempleo de los municipios del Cauca y del Valle del Cauca con estimación de áreas pequeñas, integrando varias fuentes y covariables con respaldo en la literatura. Se desarrolló y se aplicó al año 2018: un procedimiento de 16 etapas que integra la GEIH, TerriData, la DIVIPOLA y las cifras publicadas por el DANE, y que entrega una estimación para cada uno de los 83 municipios objetivo con su error, su estado de extrapolación y una calidad reportada.

El procedimiento reproduce las cifras oficiales donde existen. Las tasas directas de los 23 dominios coinciden con las del DANE, con una diferencia máxima de 0.0005 puntos porcentuales, y la tasa conjunta, 10.83~\%, es la del total de 23 ciudades y áreas metropolitanas. Los 23 dominios quedan como confiables, con CV entre 2.01~\% y 3.23~\%.

Las estimaciones municipales, en cambio, tienen un alcance limitado. Ninguno de los 83 municipios queda como confiable: 46 son aceptables y 37 no confiables por su CV, que va de 18.07~\% a 30.03~\%. De ellos, 21 son extrapolación del modelo, y todos dependen del supuesto de que la relación estimada en 23 dominios urbanos vale para municipios con una población urbana mediana de 45.07~\%. El aporte del trabajo es, por eso, la metodología y no las cifras: un procedimiento reproducible que dice con qué calidad y bajo qué supuestos se puede estimar el desempleo de un municipio sin encuesta, y que puede repetirse con mejores covariables u otros años.

%--------------------------------------------------------------------------------------
\section{Conclusión asociada al objetivo específico 1}

El primer objetivo específico pedía evaluar el desempeño de modelos de áreas pequeñas con componentes espaciales. Se compararon el Fay-Herriot clásico, elegido entre cuatro especificaciones candidatas, y ocho especificaciones espaciales SAR, con cuatro matrices de proximidad por distancia y dos estructuras. Todas usan la misma respuesta, los mismos $\psi_d$, los mismos 23 dominios y la misma covariable.

Los datos no dieron evidencia de dependencia espacial: ninguna especificación tiene un LRT de $\rho = 0$ significativo ($p \geq 0.221$), el I de Moran de los residuos no es significativo ni en el clásico (I de 0.093 con 3 vecinos, $p = 0.333$) ni en el espacial, y el AIC del espacial es mayor en todos los casos. En los dominios las dos estimaciones son prácticamente iguales (diferencia máxima de 0.02 puntos porcentuales y CV mediano de 2.60~\% en ambos); en los municipios la diferencia absoluta media es 0.22 puntos. El modelo recomendado es el Fay-Herriot clásico con la tasa de homicidios, y el espacial queda como análisis de sensibilidad. Con 23 ciudades dispersas la prueba tiene poca potencia, de modo que el resultado dice que estos datos no justifican el componente espacial, no que el desempleo municipal carezca de estructura espacial.

%--------------------------------------------------------------------------------------
\section{Conclusión asociada al objetivo específico 2}

El segundo objetivo específico pedía determinar las variables con mayor incidencia en la tasa de desempleo. De 18 candidatas con respaldo en la literatura, la única que quedó en el modelo es la tasa de homicidios, con una asociación positiva: 0.099 puntos porcentuales más de desempleo por cada homicidio adicional por 100,000 habitantes (error estándar de 0.036), que se mantiene casi igual en el modelo espacial (0.102). Es una asociación estadística condicional a la selección, no un efecto causal.

Esa asociación es frágil. Explica el 27~\% de la variación entre dominios, depende de un solo dominio (sin Quibdó baja un 38~\% y deja de ser significativa) y no se sostiene cuando la selección se repite sin cada dominio: el procedimiento completo produce seis conjuntos distintos y estima peor que el modelo de solo intercepto (RMSE de 3.43 frente a 2.77 puntos porcentuales). La conclusión es que las covariables municipales disponibles en TerriData explican poco de las diferencias de desempleo entre ciudades, y que mejorar las covariables es la condición para que el modelo aporte más que el promedio de los dominios.

%--------------------------------------------------------------------------------------
\section{Conclusión asociada al objetivo específico 3}

El tercer objetivo específico pedía analizar la distribución y los patrones espaciales del desempleo municipal a partir de las estimaciones. Las 83 estimaciones van de 9.19~\% a 31.17~\%, con mediana de 13.29~\%; el Cauca y el Valle del Cauca tienen promedios casi iguales (14.16~\% y 14.18~\%), pero el Cauca es más disperso y tiene más municipios que extrapolan (13 de 41 frente a 8 de 42). Las tasas más altas se concentran en el norte del Cauca y en el norte del Valle, todas en municipios que extrapolan.

Las estimaciones tienen autocorrelación espacial positiva (I de Moran de 0.265, $p = 0.001$), pero ese patrón es el de la tasa de homicidios, que con la misma matriz da 0.256, y el efecto espacial del modelo es menor que 0.5 puntos porcentuales en 75 de los 83 municipios. El mapa de las estimaciones describe, por tanto, la distribución territorial de la covariable usada y no una medición del desempleo en esos municipios; su lectura como patrón del desempleo queda sujeta a que la relación del modelo valga fuera de los dominios.
